
\section{符号说明}

本文使用的主要符号及其含义如表\ref{tab:符号说明}所示。

\begin{longtable}{>{\centering\arraybackslash}p{0.15\textwidth}>{\centering\arraybackslash}p{0.79\textwidth}}
  \caption{符号说明}
  \label{tab:符号说明} \\
  \toprule
  符号 & 说明 \\
  \midrule
  \endfirsthead
  \caption{符号说明（续）} \\
  \toprule
  符号 & 说明 \\
  \midrule
  \endhead
  \bottomrule
  \endfoot
  $L$ & 微构体边长，$L=10000$ nm \\
  $R_A$ & 金属A圆柱体底面半径，$R_A=30$ nm \\
  $L_A$ & 金属A圆柱体长度，$L_A=5000$ nm \\
  $R_B$ & 金属B球体半径，$R_B=200$ nm \\
  $d_t$ & 隧穿距离，$d_t=1.8$ nm \\
  $N_A$ & 金属A圆柱体数量 \\
  $N_B$ & 金属B球体数量 \\
  $V_A$ & 单个金属A的体积，$V_A=\pi R_A^2 L_A$ \\
  $V_B$ & 单个金属B的体积，$V_B=\frac{4}{3}\pi R_B^3$ \\
  $\phi$ & 填充率（体积分数），$\phi = (N_A V_A + N_B V_B) / L^3$ \\
  $P_c$ & 导通概率 \\
  $\mathbf{p}_1, \mathbf{p}_2$ & 圆柱体两端点三维坐标 \\
  $d_{AA}$ & 两圆柱体轴线间最短距离 \\
  $d_{AB}$ & 球心到圆柱体轴线的最短距离 \\
  $d_{BB}$ & 两球心间最短距离（考虑周期边界） \\
  $d_{\text{surf}}$ & 两颗粒表面间最短距离 \\
  $\mathcal{G}$ & 微构体区域 $[-5000,5000]^3$ nm$^3$ \\
  $M$ & 蒙特卡洛模拟次数 \\
  $\hat{P}_c$ & 导通概率的蒙特卡洛估计值 \\
\end{longtable}
